\section{Bounded nonmyopic probing and its limits}
\label{sec:v6}
Version 6 asks whether two-observation lookahead repairs the greedy plateau at
useful cost. This phase uses finite CPU simulators with every neural ranker
frozen; visit counters and exact evidence still update after actual episodes.
No new language-model inference, weight adaptation, deployment, or scheduled
continuation is part of this phase.

\paragraph{Established antecedents.}
Observation-contingent policy trees and belief-state planning are established
POMDP methods~\cite{pomdp}. POMCP and Bayes-adaptive sampled search already value
information through future decisions while avoiding full tree
expansion~\cite{pomcp,bamcp}. Their sampled objectives do not certify a worst-case
bound over every compatible model. Adaptive submodularity provides conditions
for effective greedy acquisition; the parity counterexample violates the
relevant diminishing-returns intuition~\cite{adaptivesubmod}. EC$^2$ instead cuts
edges between hypotheses that require different decisions, so useful tests need
not immediately improve attainable reward~\cite{ec2}. Overlapping acceptable
decision regions require additional care, as in hyperedge cutting~\cite{hec}.
Neither lookahead nor decision-directed acquisition is a novelty claim here.

\paragraph{Algorithm and retained authority.}
All three v6 selectors first attempt an exact strict or weak improvement using
the unchanged paired checker. When that fails, the depth-two selector enumerates
the compatible models and plans a first probe with a different second probe
allowed for each observed trace. For the library regret $R(C)$ defined in the v5
study, a contingent plan $\pi=(p,\{q_\tau\})$ has score
\[
 S(\pi)=\frac{R(C)-\max_{\tau,\sigma}R(C_{p,\tau,q_\tau,\sigma})}
 {1+d_p+\max_\tau d_{q_\tau}}.
\]
Stopping is allowed in every branch; all debits use exact comparisons with the
current incumbent and every branch must fit the remaining budget. The current
incumbent and goal remain fixed within this hypothetical plan. For each first
probe, we enumerate attainable terminal-regret thresholds and choose the cheapest
feasible second alternative in each branch under each threshold. This finite
procedure covers the optimum of the stated gain/debit surrogate. Independently
minimizing regret in each branch can waste debit below the worst branch's
bottleneck; a counterexample and regression tests reject that earlier variant.
The whole selector is still restricted by its promotion shortcut, program
library, proposal set, and computation caps.

Only the first action program executes, after a final exact comparison. The
second-probe map is an explanation, not a future execution ticket. Real feedback
causes replanning, so later goal changes or computation limits can prevent the
predicted two-step gain from materializing. Simulated traces never enter the
experience record. A plan that exhausts its computation allowance returns
\unk, keeps the incumbent, and spends zero additional debit.

The cheap control chooses an informative program with a nonnegative exact lower
return difference, including reward ties. It uses possible-trace search without
enumerating complete models. A second control maximizes optimistic current
return under the same exact checks and B16 ledger. All three have the same finite
inputs and a common planning ceiling: two million package Python line events,
a cooperative one-second deadline, 4,096 distinct models, and 100,000 nodes per
paired check. Counters include attempted duplicate completions and interrupted
search work. They exclude their own bookkeeping, construction, observation
updates, and C-backed/library internals. This is neither a FLOP metric nor an
operating-system resource sandbox. Actual elapsed costs are reported separately.

\paragraph{Positive and negative controls.}
Every selector obtains returns $(5,5,10)$ on all four two-bit parity worlds at
B0. Thus two-step planning repairs the motivating failure, but deeper planning
is unnecessary for that repair. In the analogous three-bit class, every B0-feasible plan of at most two
probes has zero regret reduction. The depth-two rule remains at return
five in all eight worlds, while the cheap rule obtains $(5,5,5,10)$. The same
construction defeats every fixed planning depth at B0 when it contains more
hidden bits than the lookahead. Equal-payoff controls also charge extra planning
without any possible reward gain. These are declared development mechanisms,
not independent benchmark validation.

\paragraph{Frozen held-out protocol.}
The primary family contains all 1,296 labelled two-state, two-action,
three-output transducers. The library has eight length-three action words, and
the public objectives have rewards $(0,1,2)$ and $(2,1,0)$. Each 24-episode stream
uses eight episodes of A, eight of B, then eight of A; future schedule and true
world tables never enter the learner. Development uses seeds 60000--60007;
40 held-out streams use seeds 61000--61039 after a local source/configuration
freeze. Seeds sample tables with replacement. Correcting the meter's path
handling and the branch-debit scoring defect occurred before this freeze;
the initial and corrected development outputs are both retained. This is local
protocol recording, not external preregistration.

\begin{table}[t]
\centering\small\setlength{\tabcolsep}{4.5pt}
\begin{tabular}{lrrrr}
\toprule
Method & Return & Plan s & Unknown & Max debt\\
\midrule
Depth two & 4.197 & 8.344 & 30.7\% & 0\\
Free info. & 3.803 & 0.940 & 0.0\% & 0\\
Bounded opt. & 4.392 & 1.560 & 0.1\% & 9\\
V4 frozen & 4.391 & 0.010 & --- & 6\\
Full history & 4.432 & 0.036 & --- & 6\\
\bottomrule
\end{tabular}
\caption{Frozen v6 holdout: mean episode return, total planning seconds per 24-episode stream, and maximum completed-prefix anchor deficit. Only the first three arms share line metering and the one-second ceiling. Full-history optimism has no risk ledger. All ranker weights stay frozen. Timing includes measurement overhead.}
\label{tab:v6}
\end{table}
The primary mean difference is $+0.39375$ reward per episode, with paired 95\% Student interval $[+0.15845,+0.62905]$ over 40 independent streams. The depth-two selector uses 8.88$\times$ the cheap control's measured planning time. The holdout contains 40 distinct tables; 0 streams match a development table (0 distinct overlapping tables). This is same-family validation, not new-domain generalization.
\begin{figure}[t]
\centering
\includegraphics[width=\columnwidth]{figures/v6_reward_cost.pdf}
\caption{Reward and measured planning cost for the three commonly metered v6 arms. Reward intervals are marginal 95\% Student intervals across streams; the primary analysis uses paired differences. Common ceilings do not imply equal realized computation.}
\label{fig:v6cost}
\end{figure}


The positive primary contrast does not establish efficiency superiority. Relative
to bounded optimism, depth two loses $0.19479$ reward per episode (paired 95\%
interval $[-0.31761,-0.07198]$); it also loses $0.19375$ to the frozen-ranker v4
control and $0.23542$ to unbounded full-history optimism. These are secondary,
unadjusted contrasts. Mean charged risk is 1.775 for depth two, zero for the free
control, and 6.975 for bounded optimism. The common budget is an upper ceiling,
so the positive contrast does not isolate lookahead from willingness to spend
risk. No guarded class, incumbent, or budget audit fails; all 3,840 guarded
holdout prefixes satisfy B16. Depth two returns \unk\ on 30.7\% of decisions.
The efficiency-superiority hypothesis is unsupported, while the narrowly stated
reward contrast against the zero-loss rule is supported on this family.

\paragraph{An executable formal bridge.}
The Lean inventory now has 63 kernel-checked declarations, including 12 new
statements about an executable finite transducer interpreter and history filter.
Let $\operatorname{run}(M,w)$ be the observable trace from reset and let
$H_M(W)$ be the history generated by actually running each word in $W$.
The checked characterization is
\[
\begin{aligned}
&m\in\operatorname{filterHistory}(C,H)\\
&\quad\iff m\in C\;\land\;\forall t\in H,\\
&\qquad\operatorname{run}(m,\operatorname{actions}(t))=t.
\end{aligned}
\]
It follows constructively that $M\in C$ implies
$M\in\operatorname{filterHistory}(C,H_M(W))$. This relates evidence retention
to execution rather than assuming an uninterpreted consistency predicate.
All 63 declarations use only the stated standard Lean dependencies and no
placeholders or custom axioms; 24 are axiom-free.

An independently compiled Lean executable supplies 7,936 traces and 1,024
history survivor sets over all 256 binary two-state machines. Python matches
all traces, all 7,680 nonempty word-program executions, and 262,144 filter
membership decisions. These are finite differential checks of the Python
bridge, not a general refinement proof of the partial-table DFS or adaptive
policy checker. The theorem still requires genuine feedback, reliable reset,
stationary deterministic dynamics, and initial inclusion of the true machine.
It establishes neither alignment nor a native continual-learning win.
